Skin cancer, and melanoma in particular, is highly treatable when detected
early, but the chances of survival deteriorate at advanced stages; at the same
time, access to specialist dermatological care is unevenly distributed worldwide.
The aim of this thesis is therefore to present a deep-learning-based
artificial-intelligence model capable of image-based classification of skin
lesions, together with the ``skin-cancer screening'' system built around it,
accessible both in a web browser and on mobile devices, which indicates the extent
to which a specialist examination of a skin lesion is warranted.

As the classifier we used a Vision Transformer (ViT) model fine-tuned with
transfer learning, trained on the ISIC 2024 database, which made it suitable for
classifying non-dermoscopic skin-lesion images. We compared the model with a
ResNet-18 convolutional baseline trained under identical conditions, with the
ISIC competition-winning EVA02 model, and with the Gemini 2.5 Flash large
language model. On the test set the fine-tuned ViT reached an AUC of $0.939$,
substantially outperforming the convolutional baseline and matching the level of
EVA02.

We extended the trained classifier model into a complete application, comprising
mobile and web clients, an MCP layer serving the model, Gemini-based image
validation and natural-language explanation, and authenticated, per-user cloud
storage.

\textbf{Keywords}: skin cancer, melanoma, deep learning, Vision Transformer,
transfer learning, ISIC 2024, large language model, mobile health
